% SPDX-License-Identifier: CC-BY-SA-4.0
% Author: Matthieu Perrin

\begingroup

\SetKwFunction{Propose}{propose}
\SetKwFunction{Rand}{random\_boolean}
\SetKwData{AAC}{aac}
\SetKwVariable{Val}{val}
\newcommand\Adopt{\textup{\textsc{adopt}}}
\newcommand\Abort{\textup{\textsc{abort}}}
\newcommand\Commit{\textup{\textsc{commit}}}

\begin{exercice}[Consensus binaire randomisé\footnote{[B83] Michael Ben-Or. \emph{Another advantage of free choice for completely asynchronous agreement protocols.} PODC 1983}]

  On suppose désormais disposer d'une séquence infinie d'objets Abort-Adopt-Commit.  
  On passe maintenant à l'algorithme de Ben-Or, qui implémente le consensus binaire à partir d'objets Abort-Adopt-Commit.
  La fonction $\Rand()$ retourne $\False$ ou $\True$ avec probabilité $1/2$,
  indépendamment des autres appels à $\Rand$.

  \begin{algorithm}
    \lSVariables{}{
      $\AAC[1, 2, ...]$\tcp*[f]{Séquence infinie d'objets Abort-Adopt-Commit}
    }
    \lLVariables{\At $p_i$}{
      $\Val_i \leftarrow \False$\tcp*[f]{Valeur courante soutenue par $p_i$}
    }
    \Operation{$\Propose(v)$ \At $p_i$}{
      \nl $\Val_i\leftarrow v$\;
      \nl \For{$r = 1, 2, ...$}{
        \nl \Switch{$\AAC[r].\Propose(\Val_i)$}{
          \nl \lCase{$\Commit(v)$}{\Return $v$}
          \nl \lCase{$\Adopt(v)$~~~}{$\Val_i\leftarrow v$}
          \nl \lCase{$\Abort$~~~~~~~~}{$\Val_i\leftarrow \Rand()$}
        }
      }
    }
    \caption{Algorithme de Ben-Or}
    \label{algo:benOr}
  \end{algorithm}

  \begin{questions}

  \item Supposons qu'à la ronde $r$, un processus obtienne $\Commit(v)$.
    Quels résultats les autres processus peuvent-ils obtenir de $\AAC[r]$ ?
    Quelles valeurs les processus qui poursuivent l'exécution peuvent-ils alors proposer au tour $r+1$ ?
    En déduire qu'aucune valeur différente de $v$ ne pourra être décidée.

  \item Supposons que tous les processus commencent un tour avec la même valeur $v$.
    Que se passe-t-il pendant ce tour ?
    En déduire la propriété de validité de l'algorithme.

  \item Remplaçons provisoirement $\Rand()$ par une valeur déterministe fixe
    $b \in \{\False, \True\}$.
    Construire une exécution dans laquelle les processus ne décident jamais.
    Quel rôle joue donc le choix aléatoire ?

  \item Supposons qu'aucun processus ne décide au tour $r$.
    Montrer qu'avec une probabilité au moins $2^{-(n-1)}$, tous les processus
    corrects qui poursuivent l'exécution proposent la même valeur au tour
    $r+1$.

  \item En déduire que tout processus correct décide avec probabilité $1$.
    Donner un majorant asymptotique du nombre moyen de tours nécessaires
    avant la décision.

  \end{questions}

\end{exercice}

\endgroup
\endinput
